\documentclass[11pt,a4paper]{article}

%% PACHETE MATEMATICE
\usepackage{amsmath,amssymb,amsfonts}
\usepackage{amsthm}
\usepackage{mathpazo}
\usepackage{eulervm}
\usepackage{enumerate}
\usepackage{color}
\usepackage{graphicx}
\usepackage[all]{xy}
\usepackage{xcolor}
\usepackage{framed}
\usepackage[unicode]{hyperref}
\setcounter{MaxMatrixCols}{20}
\usepackage{multicol}
% \usepackage[romanian]{babel}


%% ASPECT
\linespread{1.05}
\usepackage[T1]{fontenc}
\usepackage[utf8x]{inputenc}
\usepackage[margin=2cm]{geometry}
% \usepackage[color]{showkeys}
\usepackage{anyfontsize}
\usepackage{titlesec}
\titleformat*{\section}{\large\bfseries}

\usepackage{fancyhdr}
\pagestyle{fancy}
\rhead{\small \color{gray} Notițe de Adrian Manea}
\lhead{\small \color{gray} IntProb-FIA, 2017---2018, L. Lemnete \& M. Olteanu}


\usepackage[autostyle = false, style = german]{csquotes}
\usepackage[romanian]{babel}
\newcommand\qq{\enquote}

%% COMENZILE MELE
\renewcommand*{\proofname}{Dem.:}
\newcommand\bb{\mathbb}
\newcommand\dr{\mathrm}
\newcommand\kal{\mathcal}
\newcommand\fk{\mathfrak}
\newcommand\op{\oplus}
\newcommand\ot{\otimes}
\newcommand\ds{\displaystyle}
\newcommand\id{\indent}
\newcommand\nid{\noindent}
\newcommand\seq{\subseteq}
\newcommand\sneq{\subsetneq}
\newcommand\speq{\supseteq}
\newcommand\spneq{\supsetneq}
\newcommand\rar{\rightarrow}
\newcommand\Rar{\Rightarrow}
\newcommand\xrar{\xrightarrow}
\newcommand\lar{\leftarrow}
\newcommand\Lar{\Leftarrow}
\newcommand\xlar{\xleftarrow}
\newcommand\lrar{\leftrightarrow}
\newcommand\Lrar{\Leftrightarrow}
\newcommand\ideal{\trianglelefteq}
\newcommand\ai{\text{ a.\^{i}. }}
\newcommand\sk{(\textit{Schi\c{t}\u{a}}) }
\newcommand\rhup{\rightharpoonup}
\newcommand\rhdn{\rightharpoondown}
\newcommand\lhup{\leftharpoonup}
\newcommand\lhdn{\leftharpoondown}
\newcommand\vid{\emptyset}
\newcommand\wh{\widehat}
\newcommand\ol{\overline}
\newcommand\NN{\mathbb{N}}
\newcommand\RR{\mathbb{R}}
\newcommand\ZZ{\mathbb{Z}}
\newcommand\QQ{\mathbb{Q}}
\newcommand\CC{\mathbb{C}}

%% STILURILE MELE
\newtheoremstyle{dotless}{}{}{\itshape}{}{\bfseries}{:}{ }{}

\theoremstyle{dotless}
\newtheorem{theorem}{Teorem\u{a}}[section]
\newtheorem{proposition}{Propozi\c{t}ie}[section]
\newtheorem{lemma}{Lem\u{a}}[section]
\newtheorem{corollary}{Corolar}[section]
\newtheorem{exercise}{Exerci\c{t}iu}

\newtheoremstyle{dotlessdr}{}{}{}{}{\bfseries}{:}{ }{}
\theoremstyle{dotlessdr}
\newtheorem{example}{Exemplu}[section]
\newtheorem{definition}{Defini\c{t}ie}[section]
\newtheorem{remark}{Observa\c{t}ie}[section]




\begin{document}

\begin{center}
  {\large\textbf{Seminar 6 \\ Integrale triple}}
\end{center}

\vspace{1cm}

1. Considerăm mulțimea:
\[
  \Omega = \{ (x, y, z) \in \RR^3 \mid x^2 + \frac{y^2}{4} \leq 1, x^2 + y^2 \geq 1, x,y \geq 0, 0 \leq z \leq 5 \}.
\]

Calculați integrala $ \ds\iiint_\Omega \frac{yz}{\sqrt{x^2 + y^2}} dxdydz $ prin două metode:
\begin{enumerate}[(a)]
\item proiectînd mulțimea $ \Omega $ pe planul $ XOY $;
\item folosind coordonate cilindrice.
\end{enumerate}

\emph{Soluție:} (a) Pentru a proiecta mulțimea $ \Omega $ pe planul $ XOY $, studiem doar
ecuațiile corespunzătoare acestui plan. Astfel, avem intervalul $ z \in [0, 5] $, iar pentru
$ x $ și $ y $, obținem:
\[
  D = \{(x, y) \in \RR^2 \mid x \in [0, 1], \sqrt{1 - x^2} \leq y \leq 2 \sqrt{1 - x^2} \}.
\]

Atunci integrala se obține:
\begin{align*}
  \iiint_\Omega \frac{yz}{\sqrt{x^2 + y^2}} dxdydz &= \iint_D dxdy \int_0^5 \frac{yz}{\sqrt{x^2 + y^2}} dz \\
                                                   &= \frac{25}{2} \iint_D \frac{y}{\sqrt{x^2 + y^2}} dxdy,
\end{align*}
integrală pe care o calculăm folosind coordonate polare.

(b) Coordonatele cilindrice sînt:
\[
  \begin{cases}
    x &= \rho \cos \varphi \\
    y &= \rho \sin \varphi \\
    z & = z
  \end{cases}, \quad (\rho, \varphi, z) \in [0, \infty) \times [0, 2\pi) \times \RR,
\]
iar jacobianul transformării este $ J = \rho $.

Pentru $ \Omega $, avem:
\[
  z \in [0, 5], \quad \varphi \in \big[ 0, \frac{\pi}{2} \big], %
  \rho \in \big[ 1, \frac{2}{\sqrt{3 \cos^2 \varphi + 1}} \big].
\]

Obținem, atunci:
\begin{align*}
  \iiint_\Omega \frac{yz}{\sqrt{x^2 + y^2}} dxdydz &= \int_0^5 dz \int_0^{\frac{\pi}{2}} \int_1^{\frac{2}{\sqrt{3\cos^2 \varphi + 1}}} %
                                                     \frac{z \rho \sin \varphi}{\rho} \rho d\rho \\
                                                   &= \frac{25}{2} \int_0^{\frac{\pi}{2}} \frac{3(1 - \cos^2 \varphi)}{3 \cos^2 \varphi + 1} \sin \varphi d\varphi \\
                                                   &= \frac{25}{18} (4 \pi \sqrt{3} - 9).
\end{align*}

\vspace{1cm}

2. Folosind coordonatele sferice (generalizate) să se calculeze integralele:
\begin{enumerate}[(a)]
\item $ \ds\iiint_\Delta ( y - x ) dxdydz $, unde domeniul este:
  \[
    \Delta = \{(x, y, z) \in \RR^3 \mid x^2 + y^2 + z^2 \leq 1, y > 0 \};
  \]
\item $ \ds\iiint_\Omega \Big( 1 - x^2 - \frac{y^2}{9} - \frac{z^2}{4} \Big)^{\frac{3}{2}} dxdydz $,
  domeniul fiind:
  \[
    \Omega = \{ (x, y, z) \in \RR^3 \mid y \geq 0, z \geq 0, x^2 + \frac{y^2}{9} + \frac{z^2}{4} \leq 1 \};
  \]
\item $ \ds\iiint_\Pi z dxdydz $, unde domeniul este:
  \[
    \Pi = \{(x, y, z) \in \RR^3 \mid x^2 + y^2 + (z - 1)^2 \leq 1 \}.
  \]
\end{enumerate}

\emph{Indicații:} Coordonatele sferice sînt:
\[
  \begin{cases}
    x &= \rho \sin \theta \cos \varphi \\
    y &= \rho \sin \theta \sin \varphi \\
    z &= \rho \cos \theta
  \end{cases}, \quad (\rho, \theta, \varphi) \in [0, \infty) \times [0, \pi] \times [0, 2\pi).
\]

Jacobianul transformării este $ J = \rho^2 \sin \theta $.

Pentru coordonatele sferice generalizate (utile pentru elipsoizi) sînt:
\[
  \begin{cases}
    x &= a \rho \sin \theta \cos \varphi \\
    y &= b \rho \sin \theta \sin \varphi \\
    z &= c \rho \cos \theta
  \end{cases},
\]
definite pe același domeniu de mai sus, iar jacobianul este $ J = abc\rho^2 \sin \theta $.

\vspace{1cm}

3. Calculați volumele mulțimilor $ \Omega $, mărginite de suprafețele de ecuații:
\begin{enumerate}[(a)]
\item $ 2x^2 + y^2 + z^2 = 1, \quad 2x^2 + y^2 - z^2 = 0, z \geq 0 $;
\item $ z = x^2 + y^2 - 1, \quad z = 2 - x^2 - y^2 $;
\item $ z = 4 - x^2 - y^2, \quad 2z = 5 + x^2 + y^2 $;
\item $ x^2 + y^2 = 1, \quad z^2 = x^2 + y^2, \quad z \geq 0 $.
\end{enumerate}

\emph{Indicații:} (a) Intersectăm cele două suprafețe, egalînd primele două ecuații și
obținem:
\[
  4x^2 + 2y^2 = 1, \quad z = \frac{1}{\sqrt{2}}.
\]
Atunci, proiecția acestui domeniu pe planul $ XOY $ este:
\[
  D = \{(x, y) \mid 4x^2 + 2y^2 \leq 1 \}.
\]
Iar coordonata $ z $ o putem lua între cele două suprafețe, adică:
\[
  \sqrt{2x^2 + y^2} \leq z \leq \sqrt{1 - 2x^2 - y^2}.
\]
Rezultă că putem calcula volumul mulțimii astfel:
\[
  V(\Omega) = \iiint_\Omega dxdydz = \iint_D dxdy \int_{\sqrt{2x^2 + y^2}}^{\sqrt{1 - 2x^2 - y^2}} dz,
\]
iar apoi pe domeniul $ D $ putem folosi coordonate polare generalizate (pentru parametrizarea
elipsei).

\vspace{1cm}

(b) Similar, intersectăm cele două suprafețe rezolvînd sistemul de ecuații
și găsim:
\[
  x^2 + y^2 = \frac{3}{2}, \quad z = \frac{1}{2}.
\]
Proiecția pe planul $ XOY $ este:
\[
  D = \big\{ (x, y) \mid x^2 + y^2 \leq \frac{3}{2} \big\}.
\]
Atunci:
\[
  V(\Omega) = \iiint_\Omega dxdydz = \iint_D dxdy \int_{x^2 + y^2 - 1}^{2 - x^2 - y^2} dz.
\]



\end{document}